% ============================================================================
% Perspective: Why Data-Driven Modeling of Limit Cycles Is Hard
% Target venue: Philosophical Transactions of the Royal Society A
% ============================================================================
\documentclass[12pt]{article}

% --- Packages ---------------------------------------------------------------
\usepackage[margin=1in]{geometry}
\usepackage{amsmath,amssymb,amsthm}
\usepackage{mathtools}
\usepackage{bm}
\usepackage{graphicx}
\usepackage[dvipsnames]{xcolor}
\usepackage{hyperref}
\usepackage{cleveref}
\usepackage{booktabs}
\usepackage{multirow}
\usepackage{enumitem}
\usepackage{natbib}
\usepackage{caption}
\usepackage{subcaption}
\usepackage{algorithm}
\usepackage{algpseudocode}
\usepackage{tikz}
\usetikzlibrary{arrows.meta,positioning,decorations.markings}

% --- Custom commands --------------------------------------------------------
\newcommand{\R}{\mathbb{R}}
\newcommand{\C}{\mathbb{C}}
\newcommand{\N}{\mathbb{N}}
\newcommand{\T}{\mathbb{T}}
\newcommand{\norm}[1]{\lVert #1 \rVert}
\newcommand{\abs}[1]{\lvert #1 \rvert}
\newcommand{\inner}[2]{\langle #1, #2 \rangle}
\newcommand{\op}[1]{\hat{#1}}
\newcommand{\Atilde}{\tilde{A}}
\newcommand{\ie}{\textit{i.e.}\@\xspace}
\newcommand{\eg}{\textit{e.g.}\@\xspace}
\DeclareMathOperator{\diag}{diag}
\DeclareMathOperator{\svd}{svd}
\DeclareMathOperator{\rank}{rank}
\DeclareMathOperator{\spec}{spec}
\DeclareMathOperator{\opnorm}{opnorm}
\DeclareMathOperator*{\argmin}{arg\,min}
\DeclareMathOperator*{\argmax}{arg\,max}

% Colors for annotations
\newcommand{\todo}[1]{\textcolor{red}{[\textbf{TODO:} #1]}}
\newcommand{\note}[1]{\textcolor{blue}{[\textbf{Note:} #1]}}

% --- Document ---------------------------------------------------------------
\begin{document}

% ============================================================================
\title{%
  Non-normality and the Limits of Data-Driven Decomposition \\
  for Stable Oscillators: \\[4pt]
  \large A Perspective from Gait Biomechanics%
}

\author{%
  Michael Hess\textsuperscript{1,*},
  Gordon Berman\textsuperscript{1,2}
  \\[6pt]
  \textsuperscript{1}[Department of Biology, Emory University] \\
  \textsuperscript{2}[Department of Physics, Emory University] \\
  \textsuperscript{*}Corresponding author: \texttt{[michael.hess@emory.edu]}
}

\date{}

\maketitle

% ============================================================================
\begin{abstract}
Limit cycles, or stable periodic orbits, are ubiquitous in biological systems,
from cardiac rhythms and neural oscillations to the rhythmic coordination of
locomotion. Recent advancements in data-driven methods such as Dynamic Mode Decomposition (DMD) and
neural ordinary differential equations (ODEs) promise to extract mechanistic
insight from time-series observations of these systems, yet in practice they
encounter a constellation of difficulties that are rarely discussed together.
We present a unified perspective on why decomposing limit-cycle dynamics from
data is fundamentally harder than the analogous problem for fixed-point
attractors. Three interacting mechanisms conspire against standard approaches:
(i) measurement noise systematically biases identified linear operators,
shrinking eigenvalues and corrupting mode shapes in ways that total-least-squares
methods only partially remedy; (ii) the Floquet structure of limit cycles
induces \emph{intrinsic non-normality} in any single-operator approximation,
so that eigenvalue decomposition fails to separate dynamical contributions even
when the spectrum is perfectly estimated; and (iii) the phase-dependent geometry
of stable isostable foliations---twisting Floquet eigendirections and varying
contraction rates---means that no time-invariant linear model can faithfully
represent the transverse dynamics without mode mixing. We develop diagnostic
tools based on pseudospectral analysis and resolvent gain curves, and apply
them to a Hankel-DMD pipeline for human locomotion data (299 gait trials
across able-bodied controls and stroke survivors). We show that non-normality
metrics derived from the DMD operator discriminate clinical populations with
large effect sizes ($p < 0.001$), suggesting that the very structure that
makes limit cycles hard to decompose also encodes clinically relevant
information about gait stability. We discuss implications for neural ODE
approaches, spectral submanifold methods, and piecewise-linear strategies,
arguing that the community should embrace non-normality as a feature rather
than a nuisance.

\vspace{6pt}
\noindent\textbf{Keywords:}
Dynamic Mode Decomposition, limit cycles, non-normality, Floquet theory,
pseudospectra, gait biomechanics, stroke rehabilitation, neural ODEs
\end{abstract}

% ============================================================================
\section{Introduction}
\label{sec:intro}

% --- Opening: limit cycles are everywhere ---
Rhythmic behavior is among the most fundamental modes of biological
organization. The heartbeat, respiratory rhythm, circadian clocks, neural
oscillations, and locomotor coordination are all examples of \emph{limit
cycles}---isolated periodic orbits that attract nearby trajectories
\citep{strogatz2015nonlinear, guckenheimer1983nonlinear}. Unlike fixed points,
limit cycles involve continuous motion around a closed orbit, with perturbations
relaxing back to the cycle through transverse contracting directions described
by \emph{Floquet theory} \citep{floquet1883equations, chicone2006ordinary}.
Understanding the geometry of this relaxation---the Floquet multipliers
and their associated eigendirections---is critical for predicting how
biological oscillators respond to perturbation, for designing therapeutic
interventions, and for building predictive models.

% --- The data-driven promise ---
The past two decades have seen an explosion of data-driven methods for
dynamical systems. Dynamic Mode Decomposition (DMD)
\citep{schmid2010dynamic, kutz2016dynamic} extracts a best-fit linear operator
from time-series snapshots; its variants (extended DMD, Hankel-DMD, optimized
DMD) have been applied to fluid flows, neuroscience, epidemiology, and
robotics. Neural ordinary differential equations (neural ODEs)
\citep{chen2018neural} offer a nonlinear alternative, learning continuous-time
dynamics from data through backpropagation. Spectral submanifold (SSM) methods
\citep{haller2016nonlinear, cenedese2022data} seek to identify the
lowest-dimensional invariant manifolds that capture the essential dynamics.
All three families of methods have been motivated, benchmarked, and primarily
validated on systems with \emph{fixed-point attractors}---damped oscillations
returning to equilibrium.

% --- The hard truth ---
When these methods are applied to \emph{limit cycles}, they encounter a
qualitatively different landscape of difficulties that the fixed-point
intuition does not prepare practitioners for. The purpose of this perspective
is to lay out these difficulties in a unified framework, grounding the
discussion in both theory and a concrete application to human gait data.
Our central thesis is:

\begin{quote}
\emph{The combination of noise bias, intrinsic non-normality, and
phase-dependent transverse geometry makes limit cycles fundamentally
resistant to data-driven decomposition into separable, interpretable
components---and this resistance is itself a meaningful dynamical signature
that can be exploited clinically.}
\end{quote}

% --- Roadmap ---
We organize the perspective as follows.
\Cref{sec:background} reviews the essential theory: Floquet analysis,
DMD and its noise sensitivity, and the sources of non-normality in
single-operator limit-cycle models.
\Cref{sec:methods} describes our Hankel-DMD pipeline with stability
enforcement, and the pseudospectral and resolvent-based diagnostic toolkit
we have developed.
\Cref{sec:results_sl} presents a controlled study on the Stuart--Landau
oscillator, where ground-truth Floquet structure is known, demonstrating
how non-normality emerges even from a nearly diagonal system.
\Cref{sec:results_gait} applies the full pipeline to 299 gait trials
from able-bodied controls and stroke survivors, showing that non-normality
metrics discriminate clinical populations.
\Cref{sec:alternatives} discusses why neural ODEs, piecewise-linear
models, and spectral submanifold methods face related or compounded
difficulties.
\Cref{sec:discussion} synthesizes the results and argues for a shift
in how the community approaches data-driven limit-cycle modeling.

% --- Prior work on gait ---
This work builds on Winner \emph{et al.} \citep{winner2023lstm}, who used
LSTMs to predict and find low-dimensional structure in human locomotion
joint-angle data, demonstrating that data-driven methods can capture gait
dynamics but leaving open the question of \emph{mechanistic interpretability}
---what do the learned representations correspond to in dynamical systems
terms?

% ============================================================================
\section{Background: Why Limit Cycles Are Different}
\label{sec:background}

% ----------------------------------------------------------------------------
\subsection{Floquet Theory and the Monodromy Matrix}
\label{sec:floquet}

Consider a dynamical system
\begin{equation}
  \dot{\bm{x}} = \bm{f}(\bm{x}), \qquad \bm{x} \in \R^n,
  \label{eq:ode}
\end{equation}
possessing a stable limit cycle $\gamma$ with period $T$. Linearizing about
a point $\bm{x}_0 \in \gamma$ yields the variational equation
\begin{equation}
  \dot{\bm{\xi}} = \bm{J}(\bm{x}_0(t))\,\bm{\xi},
  \label{eq:variational}
\end{equation}
where $\bm{J}(\bm{x}_0(t)) = D\bm{f}|_{\bm{x}_0(t)}$ is the
\emph{time-periodic} Jacobian. By Floquet's theorem, the fundamental matrix
solution $\Phi(t)$ can be factored as
\begin{equation}
  \Phi(t) = P(t) \, e^{Bt},
  \label{eq:floquet}
\end{equation}
where $P(t) = P(t + T)$ is $T$-periodic and $B$ is a constant matrix.
The \emph{monodromy matrix} $M = \Phi(T)$ maps perturbations through one
full period; its eigenvalues $\mu_i$ are the \emph{Floquet multipliers},
and the corresponding eigenvectors define the directions along which
perturbations grow or decay.

One multiplier is always $\mu_1 = 1$ (the marginal direction tangent to
the orbit). The remaining multipliers $\abs{\mu_i} < 1$ for a stable
limit cycle, and their associated \emph{isostable directions}
\citep{wilson2016isostable, kvalheim2021existence} define the transverse
contracting geometry. Crucially:

\begin{enumerate}[label=(\roman*)]
  \item The Floquet eigendirections are generally \emph{not constant}---they
  rotate and stretch as one moves along $\gamma$, parameterized by the
  periodic matrix $P(t)$.
  \item The contraction rates can be \emph{phase-dependent}: a perturbation
  may transiently grow in certain phases of the cycle before ultimately
  decaying.
  \item When multipliers are close in magnitude (small separation of
  timescales), the corresponding eigendirections can become nearly
  parallel, inducing geometric sensitivity.
\end{enumerate}

These features have no analogue in fixed-point dynamics, where the
eigenvectors of the Jacobian are constant and the contraction is uniform
in time.

% ----------------------------------------------------------------------------
\subsection{DMD as a Single-Operator Approximation}
\label{sec:dmd_background}

Standard DMD \citep{schmid2010dynamic} seeks a matrix $\Atilde$ such that
$\bm{x}_{k+1} \approx \Atilde \bm{x}_k$. Given snapshot matrices
$X = [\bm{x}_1 \cdots \bm{x}_{m-1}]$ and $Y = [\bm{x}_2 \cdots \bm{x}_m]$,
the solution is
\begin{equation}
  \Atilde = Y X^+ \approx U_r^* Y V_r \Sigma_r^{-1},
  \label{eq:dmd}
\end{equation}
where $X = U \Sigma V^*$ is the SVD and $r$ is the truncation rank.

For a system near a fixed point, $\Atilde$ approximates the matrix
exponential $e^{J \Delta t}$, and the DMD eigenvalues and modes correspond
directly to the eigenvalues and eigenvectors of the Jacobian $J$.
The relationship is clean: the spectrum of $\Atilde$ carries the full
dynamical information.

For a limit cycle, the situation is fundamentally different. The data
$\bm{x}_k$ lie on or near a closed curve, and a single linear operator
$\Atilde$ must simultaneously:
\begin{enumerate}[label=(\alph*)]
  \item Rotate states around the orbit (the oscillatory component),
  \item Contract perturbations transverse to the orbit (the Floquet
  decay), and
  \item Approximate the \emph{time-varying} Floquet structure with a
  \emph{time-invariant} operator.
\end{enumerate}
Requirement (c) is the crux of the difficulty. The time-varying Jacobian
$\bm{J}(t)$ is being compressed into a constant matrix, and the resulting
$\Atilde$ must use \emph{non-normal mode mixing} to reproduce the
net-stable, phase-dependent dynamics.

% --- Takens embedding / Hankel DMD ---
\subsubsection{Hankel-DMD and Delay Embedding}
\label{sec:hankel_dmd}

To increase the rank and observability of the linear representation,
one typically applies Takens-style delay embedding
\citep{takens1981detecting, brunton2017chaos}, constructing a Hankel
(block-row) matrix:
\begin{equation}
  H = \begin{bmatrix}
    \bm{x}_1 & \bm{x}_2 & \cdots & \bm{x}_{m-\tau} \\
    \bm{x}_2 & \bm{x}_3 & \cdots & \bm{x}_{m-\tau+1} \\
    \vdots & & & \vdots \\
    \bm{x}_{\tau+1} & \bm{x}_{\tau+2} & \cdots & \bm{x}_m
  \end{bmatrix}
  \in \R^{n(\tau+1) \times (m-\tau)},
  \label{eq:hankel}
\end{equation}
where $\tau$ is the number of delay lags and $n$ is the ambient dimension.
DMD is then applied to consecutive column pairs of $H$. This ``Hankel-DMD''
approach dramatically expands the state space, enabling the linear operator
to represent richer dynamics---but it also amplifies the non-normality
problem, as we will show, because the delay coordinates introduce
additional geometric entanglement.

% ----------------------------------------------------------------------------
\subsection{Noise Bias: A Universal Problem, Amplified by Oscillations}
\label{sec:noise_bias}

% --- Forward model bias ---
Consider the noisy observation model $\bm{y}_k = \bm{x}_k + \bm{\eta}_k$,
where $\bm{\eta}_k$ is i.i.d.\ measurement noise. Standard DMD solves
\begin{equation}
  \Atilde = \argmin_A \norm{Y - A X}_F,
  \label{eq:dmd_ls}
\end{equation}
which is a \emph{forward} least-squares problem: noise in $X$ is treated as
exact, while noise in $Y$ is minimized. This asymmetry systematically biases
the identified eigenvalues \emph{toward the origin} (for continuous-time,
toward $-\infty$ on the real axis), an effect well-known in the
errors-in-variables literature \citep{fuller1987measurement} and studied
specifically for DMD by \citet{dawson2016characterizing, hemati2017biasing}.

For a fixed-point attractor, this bias makes estimated decay rates slightly
too fast---a quantitative error, but not a qualitative one. For a limit
cycle, the consequences are more severe:

\begin{enumerate}[label=(\roman*)]
  \item \textbf{Spectral radius shrinkage:} The Floquet multiplier at
  $\abs{\mu} = 1$ (the phase direction) is pulled inside the unit disk,
  making the orbit appear to spiral inward. The oscillation decays
  spuriously.
  \item \textbf{Frequency corruption:} Eigenvalue bias is generally not
  radially symmetric, so oscillation frequencies can shift.
  \item \textbf{Mode-shape distortion:} Since the bias depends on the
  signal-to-noise ratio of each mode, and oscillatory modes have high
  energy, the relative bias differs across modes, distorting the
  eigenvector structure.
\end{enumerate}

% --- OptDMD / TLS approaches ---
Optimized DMD (OptDMD, \citealt{askham2018variable}),
total-least-squares DMD \citep{hemati2017biasing}, and forward-backward
DMD \citep{dawson2016characterizing} partially address this by treating
noise in both $X$ and $Y$. However:

\begin{quote}
  \emph{Neural ODEs \citep{chen2018neural} and other continuous-time
  neural network approaches are equivalent to a forward model in this
  respect: they learn $\dot{\bm{x}} = f_\theta(\bm{x})$ by minimizing
  forward prediction error, treating the state trajectory as ground truth.
  They cannot account for noise in the state observations themselves, and
  thus suffer from the same systematic bias.}
\end{quote}

This is a critical and underappreciated point. The community has largely
treated neural ODEs as a ``more flexible'' alternative to DMD, but they
inherit exactly the same noise-bias structure and additionally lack the
spectral interpretability that makes DMD's bias diagnosable.

% ----------------------------------------------------------------------------
\subsection{Non-Normality: The Eigenvalue Decomposition Fails}
\label{sec:nonnormality_theory}

Even if noise bias could be perfectly corrected, the eigenvalue decomposition
of a limit-cycle DMD operator would \emph{still} not decompose the dynamics
into separable components. The reason is \emph{non-normality}.

A matrix $A$ is \emph{normal} if $A A^* = A^* A$. Normal matrices have
orthogonal eigenvectors and their spectral decomposition is well-conditioned.
For non-normal matrices, eigenvectors can become nearly parallel, and the
resolvent $(zI - A)^{-1}$ can be enormous even far from any eigenvalue.
The degree of non-normality is measured by:

\begin{itemize}
  \item The \textbf{eigenvector condition number} $\kappa(V)$, where $V$
  is the eigenvector matrix: $\kappa(V) = \norm{V} \cdot \norm{V^{-1}}$.
  \item The \textbf{Henrici departure from normality}:
  $\delta_H(A) = \sqrt{\norm{A}_F^2 - \sum_i \abs{\lambda_i}^2} / \norm{A}_F$.
  \item The \textbf{Kreiss constant} $\mathcal{K}(A)$, bounding transient
  growth: $\sup_{n \geq 0} \norm{A^n} \geq \mathcal{K}(A)$
  \citep{trefethen2005spectra}.
  \item The \textbf{$\varepsilon$-pseudospectral radius}
  $\rho_\varepsilon(A) = \max\{\abs{z} : z \in \sigma_\varepsilon(A)\}$.
\end{itemize}

\subsubsection{Why Limit-Cycle DMD Operators Are Necessarily Non-Normal}
\label{sec:why_nonnormal}

Consider the simplest limit cycle: the Stuart--Landau oscillator
\begin{equation}
  \dot{z} = (\sigma + i\omega) z - z\abs{z}^2,
  \label{eq:stuart_landau}
\end{equation}
where $z \in \C$, $\sigma > 0$ controls the contraction rate, and
$\omega$ is the natural frequency. On the limit cycle ($\abs{z} = \sqrt{\sigma}$),
the linearized transverse dynamics have Floquet multiplier
$\mu_s = e^{-2\sigma T}$ with $T = 2\pi/\omega$.

Even for this analytically tractable system, a Hankel-DMD operator $\Atilde$
exhibits:
\begin{itemize}
  \item $\kappa(V) \sim 10^5$ (compared to $\kappa(V) = 1$ for a normal matrix),
  \item Henrici number $\delta_H \approx 0.99$,
  \item Peak transient growth $\norm{\Atilde^n} \sim 10^4 \times$ the
  asymptotic bound.
\end{itemize}

The mechanism is clear: the operator must simultaneously encode
\emph{rotation} (eigenvalues on the unit circle) and \emph{contraction}
(eigenvalues inside the unit circle), but the \emph{exact cancellation}
required for stable periodic orbits forces the eigenvectors into near-parallel
configurations. The transient growth represents the operator's internal
mechanism for achieving the net rotation-plus-contraction: it amplifies
certain combinations of modes before they are canceled by others.

This non-normality is \textbf{not an artifact of noise, finite data, or
poor algorithm choice}---it is an \textbf{intrinsic consequence} of
approximating a time-varying process with a time-invariant operator.

\subsubsection{The Pseudospectral and Resolvent Perspective}
\label{sec:pseudospectra_theory}

The $\varepsilon$-pseudospectrum of $A$ is defined as
\begin{equation}
  \sigma_\varepsilon(A) = \{z \in \C : \sigma_{\min}(zI - A) < \varepsilon\}
  = \{z \in \C : \norm{(zI - A)^{-1}} > \varepsilon^{-1}\},
\end{equation}
and reveals the true sensitivity of the dynamics. For our discrete-time
setting (eigenvalues on or near the unit circle), we evaluate the
\textbf{resolvent norm} at $z = r e^{i\theta}$ with $r$ slightly greater
than 1:
\begin{equation}
  G(f; r) = \norm{(re^{i\theta} I - \Atilde)^{-1}}, \qquad
  f = \frac{\theta}{2\pi} f_s,
  \label{eq:resolvent_gain}
\end{equation}
where $f_s$ is the sampling frequency. This \emph{resolvent gain curve}
reveals which frequencies are most amplified by near-resonance with the
non-normal operator---a forcing-response perspective that is natural for
understanding how perturbations couple to the limit-cycle dynamics.

Key derived metrics include:
\begin{itemize}
  \item \textbf{Stability margin:}
  $d(\Atilde) = \min_\theta \sigma_{\min}(e^{i\theta}I - \Atilde)$,
  the smallest perturbation that destabilizes the system.
  \item \textbf{Kreiss constant:}
  $\mathcal{K}(\Atilde) = \sup_{\abs{z}>1} (\abs{z}-1)\norm{(zI-\Atilde)^{-1}}$,
  bounding worst-case transient growth.
  \item \textbf{Numerical abscissa:}
  $\omega(\Atilde) = \lambda_{\max}\!\bigl(\tfrac{1}{2}(\Atilde + \Atilde^*)\bigr)$,
  controlling one-step energy growth.
\end{itemize}

% ============================================================================
\section{Methods}
\label{sec:methods}

% ----------------------------------------------------------------------------
\subsection{Gait Data and Clinical Populations}
\label{sec:data}

We analyze a dataset of $N = 299$ overground walking trials from three
populations:
\begin{itemize}
  \item \textbf{AB} (Able-Bodied): 180 trials from 22 subjects, treadmill
  speeds 30--220 cm/s.
  \item \textbf{HF} (Higher-Functioning stroke): 65 trials from 11
  subjects, speeds 40--127 cm/s.
  \item \textbf{LF} (Lower-Functioning stroke): 54 trials from 9 subjects,
  speeds 22--60 cm/s.
\end{itemize}
Each trial consists of $T = 1500$ time steps at $f_s = 100$ Hz (15 seconds)
of $n = 6$ joint angles (bilateral hip, knee, ankle), preprocessed as
described in \citet{winner2023lstm}. \todo{Confirm exact joint set and
preprocessing.}

% ----------------------------------------------------------------------------
\subsection{Hankel-DMD Pipeline}
\label{sec:pipeline}

For each trial, we construct a block-Hankel matrix (\cref{eq:hankel}) with
$\tau = 10$ delay lags, yielding a state dimension of
$n_{\text{state}} = 6 \times 11 = 66$. We then:

\begin{enumerate}
  \item Compute the economy SVD $H_X = U \Sigma V^*$, where
  $H_X$ and $H_Y$ are the lagged and shifted Hankel blocks.
  \item Truncate at the noise floor: retain singular values
  $\sigma_i > 10^{-6} \sigma_1$.
  \item Form the reduced operator
  $\Atilde = U_r^* H_Y V_r \Sigma_r^{-1}$.
\end{enumerate}

\subsubsection{Stability Enforcement}
\label{sec:stability_enforcement}

The raw DMD operator frequently has eigenvalues outside the unit disk
(spectral radius $\rho > 1$), particularly for the stroke groups where
noise levels are higher. \todo{Quantify: 68\% of HF trials had
$\rho > 1$.} Since the underlying system is a stable limit cycle, these
unstable eigenvalues are artifacts. We enforce stability by projecting
each eigenvalue onto the unit disk:
\begin{equation}
  \lambda_i \leftarrow
  \begin{cases}
    \lambda_i & \text{if } \abs{\lambda_i} \leq 1, \\
    \lambda_i / \abs{\lambda_i} & \text{if } \abs{\lambda_i} > 1,
  \end{cases}
  \label{eq:stability_clamp}
\end{equation}
and reconstructing $\Atilde_{\text{stable}} = W \diag(\lambda_1, \ldots, \lambda_r) W^{-1}$.

% ----------------------------------------------------------------------------
\subsection{Non-Normality Metrics}
\label{sec:metrics}

For each trial's stability-enforced operator $\Atilde$, we compute:

\begin{enumerate}
  \item \textbf{Spectral metrics:} rank $r$, spectral radius $\rho$,
  leading singular value $\sigma_1$.
  \item \textbf{Non-normality metrics:} eigenvector condition number
  $\kappa(V)$, Henrici departure $\delta_H$, peak transient growth
  $\max_{n} \norm{\Atilde^n}$, Kreiss constant $\mathcal{K}$.
  \item \textbf{Resolvent metrics (at $r = 1.01$ and $r = 1.05$):}
  peak resolvent gain (excluding DC), mean gain, gain at 1~Hz and 2~Hz
  (the fundamental and first harmonic of gait).
  \item \textbf{Stability margin:} $d(\Atilde)$ and the frequency
  at which it occurs.
  \item \textbf{Normalized growth:} $\text{growth\_per\_rank} = \max_n \norm{\Atilde^n} / r$,
  controlling for the rank confound.
\end{enumerate}

Statistical comparisons use the Mann--Whitney $U$ test (nonparametric,
appropriate for unequal and non-normal sample sizes) with Cliff's delta
for effect size.

% ----------------------------------------------------------------------------
\subsection{Neural ODE Experiments}
\label{sec:neural_ode_methods}

\todo{Describe the neural ODE training setup: architecture (MLP layers,
hidden dim), training procedure (adjoint method), training/validation split,
and the specific failure modes observed. Key points:
\begin{itemize}
  \item Models converged to fitting the mean trajectory (the limit cycle)
  but failed to capture transverse dynamics.
  \item Forward simulation (generative mode) diverged rapidly---the models
  could not maintain a stable orbit.
  \item This is well-documented for periodic data in neural ODEs
  \citep{jordan2021neural, linot2023stabilized}.
  \item The fundamental issue: like standard DMD, neural ODEs treat $x(t)$
  as noise-free and optimize a forward loss, so they cannot account for
  errors-in-variables.
\end{itemize}}

% ============================================================================
\section{Results I: The Stuart--Landau Oscillator}
\label{sec:results_sl}

As a controlled testbed, we first apply our pipeline to the Stuart--Landau
oscillator (\cref{eq:stuart_landau}) with $\sigma = 0.1$ and $\omega = 2\pi$
(period $T = 1$ s). This system has a single nontrivial Floquet multiplier
$\mu_s = e^{-0.2\pi} \approx 0.533$ and constant isostable directions
(the radial direction in polar coordinates).

\subsection{Ground-Truth Floquet Structure vs.\ DMD Representation}

A Hankel-DMD operator with $\tau = 10$ and rank $r = 20$ yields:
\begin{itemize}
  \item Spectral radius $\rho(\Atilde) = 0.99989$ (bias from $\rho = 1$).
  \item Eigenvector condition number $\kappa(V) \approx 1.5 \times 10^5$.
  \item Henrici departure $\delta_H \approx 0.99$.
  \item Peak transient growth $\max_n \norm{\Atilde^n} \approx 1.4 \times 10^4$.
\end{itemize}

\begin{figure}[t]
  \centering
  \todo{Figure 1: Stuart--Landau DMD dissection.
  (a) Eigenvalue spectrum (unit circle, eigenvalues clustered near it).
  (b) Transient growth curve $\norm{\Atilde^n}$ vs $n$.
  (c) Pseudospectral contour plot showing $\varepsilon$-pseudospectra
      bulging far outside unit circle.
  (d) Resolvent gain curve at $r=1.01$, showing resonance peaks.}
  \caption{Non-normality anatomy of the Stuart--Landau DMD operator.
  Despite a spectral radius of 0.999, the operator exhibits transient
  amplification exceeding $10^4$ and pseudospectra that extend well
  beyond the unit circle. This non-normality is intrinsic to the
  single-operator representation of limit-cycle dynamics.}
  \label{fig:sl_anatomy}
\end{figure}

The key insight is that the Stuart--Landau system has the \emph{simplest
possible} Floquet structure---constant eigendirections, a single contraction
rate---and yet the Hankel-DMD operator is profoundly non-normal. The
non-normality arises not from the complexity of the underlying system but
from the \emph{representational mismatch}: a time-invariant linear operator
must use eigenvector non-orthogonality to encode the interplay between
oscillation and contraction.

\subsection{Mechanism: Cancellation and Mode Mixing}

The eigenvector matrix $V$ of $\Atilde$ reveals the mechanism. Modes
associated with eigenvalues near the unit circle (the ``oscillatory'' modes)
have eigenvectors that are nearly parallel to those of the ``decaying'' modes.
The operator achieves stable periodic orbits through \emph{constructive
interference on the limit cycle} (modes add coherently) and
\emph{destructive interference off the cycle} (transverse perturbations are
canceled by near-parallel mode pairs with slightly different decay rates).

This cancellation mechanism means that:
\begin{enumerate}
  \item Removing any single eigenvalue from the spectrum dramatically
  changes the predicted dynamics---modes are not separable.
  \item The transient growth $\norm{\Atilde^n}$ can vastly exceed the
  spectral-radius bound $\rho^n$, because the cancellation requires
  large intermediate amplitudes.
  \item Small perturbations to $\Atilde$ (from noise) can qualitatively
  change which modes dominate, making the decomposition unstable.
\end{enumerate}

% ============================================================================
\section{Results II: Human Gait Data}
\label{sec:results_gait}

% ----------------------------------------------------------------------------
\subsection{Stability-Enforced DMD: Basic Metrics}
\label{sec:results_basic}

Applying the pipeline to all 299 gait trials with stability enforcement
yields the results in \cref{tab:basic_metrics}.

\begin{table}[t]
\centering
\caption{Stability-enforced DMD metrics by group (median [IQR]).}
\label{tab:basic_metrics}
\begin{tabular}{lccc}
\toprule
\textbf{Metric} & \textbf{AB} ($n=180$) & \textbf{HF} ($n=65$) & \textbf{LF} ($n=54$) \\
\midrule
Rank $r$
  & 45 [43--64] & 37 [36--38] & 36 [36--37] \\
$\kappa(V)$
  & 73{,}000 [43K--121K] & 159{,}000 [113K--263K] & 136{,}000 [99K--214K] \\
$\delta_H$ (Henrici)
  & 0.98 [0.98--0.99] & 0.99 [0.93--1.00] & 0.91 [0.89--0.97] \\
$\max_n\norm{\Atilde^n}$
  & 21{,}085 [7.5K--36K] & 52{,}050 [37K--77K] & 46{,}961 [29K--68K] \\
Growth/rank
  & 475 [127--866] & 1{,}375 [1{,}026--2{,}029] & 1{,}294 [830--1{,}842] \\
\bottomrule
\end{tabular}
\end{table}

The most striking result is the \textbf{growth per rank} metric:
stroke survivors exhibit approximately $2.5$--$3\times$ the transient
growth per mode compared to able-bodied controls ($p < 0.001$, Cliff's
delta $> 0.6$, ``large'' effect). This persists in speed-matched
comparisons (30--60 cm/s), ruling out speed as a confound.

\begin{figure}[t]
  \centering
  \todo{Figure 2: Strip/jitter plots of key metrics by group.
  (a) growth\_per\_rank: AB vs HF vs LF.
  (b) $\kappa(V)$ by group.
  (c) Kreiss constant by group.
  (d) growth\_per\_rank vs speed, colored by group.}
  \caption{Non-normality metrics discriminate clinical populations.
  Growth per rank (a) shows the strongest separation ($p < 0.001$
  for AB vs HF and AB vs LF). The effect persists after controlling
  for walking speed (d).}
  \label{fig:gait_metrics}
\end{figure}

% ----------------------------------------------------------------------------
\subsection{Pseudospectral Analysis}
\label{sec:results_pseudo}

\todo{Present pseudospectral results once cell 27 computation completes.
Expected structure:
\begin{itemize}
  \item Kreiss constant: HF $>$ LF $>$ AB.
  \item Resolvent gain at $r=1.01$: stroke groups show larger peak gains,
  particularly near 1--2 Hz (gait fundamental and harmonic).
  \item Stability margin: AB has larger $d(\Atilde)$ (more robust to
  perturbation).
  \item Pseudospectral contour plots for representative AB, HF, LF trials
  showing differential ``bulging'' beyond unit circle.
\end{itemize}}

\begin{figure}[t]
  \centering
  \todo{Figure 3: Pseudospectral analysis.
  (a--c) $\varepsilon$-pseudospectral contours for representative AB, HF, LF trials.
  (d) Resolvent gain curves at $r = 1.01$ for the same three trials,
      with gait band (0.8--3 Hz) shaded.
  (e) Zoomed gait-band resolvent (0--5 Hz) at $r=1.01$.}
  \caption{Pseudospectral contours and resolvent gain curves reveal
  differential amplification structure across clinical populations.
  Stroke survivors' DMD operators show larger $\varepsilon$-pseudospectra
  (more ``leakage'' beyond the unit circle) and higher resolvent gains
  at gait-relevant frequencies.}
  \label{fig:pseudospectra}
\end{figure}

% ----------------------------------------------------------------------------
\subsection{Clinical Interpretation}
\label{sec:clinical}

The non-normality of the DMD operator is not merely a mathematical curiosity
---it encodes how the gait system responds to perturbation. Higher
growth-per-rank means that the limit-cycle dynamics rely on more delicate
cancellations between modes to maintain stability. In physical terms:

\begin{itemize}
  \item \textbf{Perturbation sensitivity:} A gait pattern with high
  non-normality can exhibit large transient excursions before returning
  to the periodic orbit, even if the orbit is formally stable (all
  Floquet multipliers inside the unit disk).
  \item \textbf{Robustness to model error:} Small changes in the operator
  (representing, e.g., fatigue, terrain changes, or neural control
  variability) have outsized effects on the predicted dynamics when the
  operator is highly non-normal.
  \item \textbf{Fall risk:} The transient growth magnitude directly
  bounds the worst-case perturbation amplification. If a stroke
  survivor's gait control system has $3\times$ the transient growth
  of an able-bodied walker, their maximum excursion from the gait
  cycle in response to an identical perturbation is correspondingly
  larger.
\end{itemize}

\todo{Discuss whether higher non-normality in stroke gait is a cause
(maladaptive control strategy) or consequence (reduced neural bandwidth
forcing the controller to operate closer to the stability boundary) of
increased fall risk.}

% ============================================================================
\section{Why Other Methods Face Similar or Worse Difficulties}
\label{sec:alternatives}

% ----------------------------------------------------------------------------
\subsection{Neural ODEs}
\label{sec:neural_odes}

Neural ODEs learn a vector field $\dot{\bm{x}} = f_\theta(\bm{x})$ by
minimizing
\begin{equation}
  \mathcal{L}(\theta) = \sum_k \norm{\bm{x}(t_k) - \hat{\bm{x}}(t_k;\theta)}^2,
  \label{eq:node_loss}
\end{equation}
where $\hat{\bm{x}}(t_k;\theta)$ is the solution of the neural ODE
integrated forward from an initial condition. This formulation has two
fundamental problems for limit cycles:

\begin{enumerate}[label=(\roman*)]
  \item \textbf{Noise-in-variables bias:} Like standard DMD, the loss
  \cref{eq:node_loss} treats the observed $\bm{x}(t_k)$ as exact. Noise
  in the state trajectory biases the learned dynamics in the same
  direction as for DMD (eigenvalue shrinkage / spurious damping), but
  the bias is harder to diagnose because the model is nonlinear and
  lacks a clean spectral decomposition.

  \item \textbf{Error accumulation:} Forward integration compounds
  errors at each time step. For a limit cycle, even a small systematic
  bias causes the predicted orbit to spiral inward or outward,
  amplified by the integration horizon. Our experiments confirm this:
  neural ODEs trained on gait data converge to the mean trajectory but
  diverge rapidly in generative mode.
  \todo{Include training curves and generative-mode divergence plots.}
\end{enumerate}

These issues have been documented for periodic data by
\citet{jordan2021neural} and \citet{linot2023stabilized}, who proposed
stabilization tricks (e.g., augmented training losses that penalize
divergence from the attractor). However, these fixes are
application-specific and do not address the fundamental noise-bias
problem.

% ----------------------------------------------------------------------------
\subsection{Piecewise-Linear Models}
\label{sec:piecewise}

An intuitive strategy for handling the time-varying Floquet structure is
to partition the limit cycle into segments and fit a separate linear model
in each:
\begin{equation}
  \bm{x}_{k+1} \approx A_j \bm{x}_k \quad \text{for } \bm{x}_k \in \mathcal{R}_j,
  \label{eq:piecewise}
\end{equation}
where $\{\mathcal{R}_j\}$ partitions the state space. This approach faces
compounding difficulties:

\begin{enumerate}[label=(\roman*)]
  \item \textbf{Phase estimation:} Defining the boundaries
  $\{\mathcal{R}_j\}$ requires knowing the limit cycle and assigning
  a phase to each data point---itself a nontrivial task for noisy,
  high-dimensional data.
  \item \textbf{Boundary discontinuities:} At region boundaries, the
  switch between linear models introduces artificial discontinuities
  that accumulate during forward propagation.
  \item \textbf{Data fragmentation:} Each submodel has access to only
  a fraction of the data, increasing noise sensitivity and making the
  bias problem \emph{worse} for each individual model.
  \item \textbf{No spectral coherence:} The eigenvalues of different
  $A_j$ need not be related in any simple way, making it impossible
  to extract global Floquet multipliers without additional assumptions.
\end{enumerate}

% ----------------------------------------------------------------------------
\subsection{Spectral Submanifolds (SSMs)}
\label{sec:ssm}

Spectral submanifold methods \citep{haller2016nonlinear, cenedese2022data}
identify the unique smoothest invariant manifolds tangent to spectral
subspaces of the linearization. They offer a rigorous nonlinear extension
of modal analysis---but with critical limitations for limit cycles:

\begin{enumerate}[label=(\roman*)]
  \item \textbf{Fixed-point formulation:} The mathematical foundation
  of SSMs requires a fixed point with a well-defined spectrum. Limit
  cycles have a trivially zero eigenvalue (the phase direction), which
  violates the non-resonance conditions needed for SSM uniqueness.

  \item \textbf{Resonances:} The oscillatory eigenvalues of a limit cycle
  are inherently resonant with the periodic orbit, creating exactly the
  conditions that SSM theory assumes away.

  \item \textbf{Adiabatic SSMs:} The Haller group has developed
  ``adiabatic SSMs'' that essentially partition the limit cycle
  (as in the piecewise approach) and compute SSMs at each phase
  \citep{li2024adiabatic}. As we noted for piecewise models, this
  reinvents the Floquet frame without resolving the fundamental
  difficulties of phase estimation and boundary management.

  \item \textbf{SINDy-based nonlinear Floquet:} Similarly, approaches
  using SINDy \citep{brunton2016discovering} to learn a nonlinear
  Floquet decomposition \citep{kvalheim2021existence} face the same
  challenge: they require accurate phase estimation and yield models
  whose coefficients vary along the orbit, deferring rather than
  resolving the decomposition problem.
\end{enumerate}

% ============================================================================
\section{Discussion}
\label{sec:discussion}

% ----------------------------------------------------------------------------
\subsection{Embracing Non-Normality}

The results of this study suggest a reframing: rather than viewing
non-normality as a failure of the DMD approximation, we should view it
as a \emph{faithful encoding} of a genuine dynamical property. A limit
cycle's transverse relaxation is inherently non-normal when viewed through
any single linear operator, and the degree of non-normality carries
information about:

\begin{enumerate}[label=(\roman*)]
  \item The separation of timescales between oscillation and contraction
  (small separation $\Rightarrow$ more non-normality).
  \item The phase-dependence and twisting of Floquet eigendirections
  (more complex geometry $\Rightarrow$ more non-normality in a
  time-invariant representation).
  \item The coupling between modes that the system relies on for
  stability (more cancellation $\Rightarrow$ higher transient growth).
\end{enumerate}

The clinical result---that stroke survivors show $\sim 3\times$ the
non-normality per mode---suggests that pathological gait is characterized
by a more delicate balance between the oscillatory and contracting
components of the dynamics. This may reflect:

\begin{itemize}
  \item Reduced neural bandwidth for real-time gait adjustments, forcing
  the system to rely on fewer, more tightly coupled modes.
  \item Asymmetric limb dynamics (hemiparesis) that rotate the Floquet
  eigendirections more dramatically around the gait cycle.
  \item Compensatory strategies that trade off robustness for energy
  efficiency.
\end{itemize}

% ----------------------------------------------------------------------------
\subsection{Implications for Data-Driven Modeling}

Our analysis suggests several practical recommendations:

\begin{enumerate}
  \item \textbf{Use total-least-squares or optimized DMD} to mitigate
  noise bias, and report the correction magnitude as a diagnostic.
  \item \textbf{Report non-normality metrics} alongside eigenvalues
  whenever applying DMD to oscillatory data. Eigenvalues alone are
  misleading for non-normal operators.
  \item \textbf{Interpret resolvent gains rather than eigenvalues} for
  frequency-domain analysis of limit-cycle systems.
  \item \textbf{Be skeptical of forward-propagation accuracy} for any
  model (DMD, neural ODE, or otherwise) applied to limit cycles.
  The transient growth bound means that errors are amplified by
  $\max_n \norm{\Atilde^n}$ before settling to the spectral decay rate.
  \item \textbf{Consider non-normality as a clinical biomarker} for
  movement disorders, complementing existing kinematic and kinetic
  measures.
\end{enumerate}

% ----------------------------------------------------------------------------
\subsection{Limitations and Future Directions}

\todo{
\begin{itemize}
  \item Subject-level analysis (current analysis is trial-level).
  \item Extension to other biological oscillators (cardiac, neural).
  \item Development of non-normality-aware DMD variants.
  \item Connection to control-theoretic robustness margins for gait.
  \item OptDMD comparison.
  \item Comparison with phase-amplitude reduction methods.
\end{itemize}}

% ============================================================================
\section{Conclusion}
\label{sec:conclusion}

Data-driven decomposition of limit-cycle dynamics is fundamentally harder
than the analogous problem for fixed-point attractors. Noise bias,
intrinsic non-normality, and phase-dependent Floquet geometry interact to
make eigenvalue-based decomposition unreliable and forward propagation
unstable. These difficulties are not specific to DMD---they affect neural
ODEs, piecewise-linear models, and spectral submanifold methods in
analogous ways. We have shown that the pseudospectral toolkit provides a
principled diagnostic framework, and that the non-normality itself encodes
clinically relevant information: stroke survivors' gait dynamics exhibit
significantly elevated transient growth, Kreiss constants, and resolvent
gains compared to able-bodied controls. We argue that the community should
shift from treating non-normality as a nuisance to exploiting it as a
window into the mechanisms of biological oscillator stability.

% ============================================================================
% References
% ============================================================================
\bibliographystyle{abbrvnat}

\begin{thebibliography}{99}

\bibitem[Askham \& Kutz(2018)]{askham2018variable}
Askham, T. \& Kutz, J.~N. (2018).
Variable projection methods for an optimized dynamic mode decomposition.
\textit{SIAM J.\ Appl.\ Dyn.\ Syst.}, 17(1), 380--416.

\bibitem[Brunton et~al.(2016)]{brunton2016discovering}
Brunton, S.~L., Proctor, J.~L., \& Kutz, J.~N. (2016).
Discovering governing equations from data by sparse identification of
nonlinear dynamical systems.
\textit{Proc.\ Natl.\ Acad.\ Sci.}, 113(15), 3932--3937.

\bibitem[Brunton et~al.(2017)]{brunton2017chaos}
Brunton, S.~L., Brunton, B.~W., Proctor, J.~L., Kaiser, E., \& Kutz, J.~N.
(2017). Chaos as an intermittently forced linear system.
\textit{Nat.\ Commun.}, 8, 19.

\bibitem[Cenedese et~al.(2022)]{cenedese2022data}
Cenedese, M., Ax{\aa}s, J., Balb{\'a}s, B., Georgiades, F., \& Haller, G.
(2022). Data-driven modeling and prediction of non-linearizable dynamics
via spectral submanifolds.
\textit{Nat.\ Commun.}, 13, 872.

\bibitem[Chen et~al.(2018)]{chen2018neural}
Chen, R.~T.~Q., Rubanova, Y., Bettencourt, J., \& Duvenaud, D. (2018).
Neural ordinary differential equations.
\textit{Advances in Neural Information Processing Systems}, 31.

\bibitem[Chicone(2006)]{chicone2006ordinary}
Chicone, C. (2006).
\textit{Ordinary Differential Equations with Applications}.
Springer.

\bibitem[Dawson et~al.(2016)]{dawson2016characterizing}
Dawson, S.~T.~M., Hemati, M.~S., Williams, M.~O., \& Rowley, C.~W. (2016).
Characterizing and correcting for the effect of sensor noise in the dynamic
mode decomposition.
\textit{Exp.\ Fluids}, 57, 42.

\bibitem[Floquet(1883)]{floquet1883equations}
Floquet, G. (1883).
Sur les \'equations diff\'erentielles lin\'eaires \`a coefficients
p\'eriodiques.
\textit{Ann.\ Sci.\ \'Ecole Norm.\ Sup.}, 12, 47--88.

\bibitem[Fuller(1987)]{fuller1987measurement}
Fuller, W.~A. (1987).
\textit{Measurement Error Models}.
Wiley.

\bibitem[Guckenheimer \& Holmes(1983)]{guckenheimer1983nonlinear}
Guckenheimer, J. \& Holmes, P. (1983).
\textit{Nonlinear Oscillations, Dynamical Systems, and Bifurcations of
Vector Fields}.
Springer.

\bibitem[Haller \& Ponsioen(2016)]{haller2016nonlinear}
Haller, G. \& Ponsioen, S. (2016).
Nonlinear normal modes and spectral submanifolds: Existence, uniqueness
and use in model reduction.
\textit{Nonlinear Dyn.}, 86(3), 1493--1534.

\bibitem[Hemati et~al.(2017)]{hemati2017biasing}
Hemati, M.~S., Rowley, C.~W., Deem, E.~A., \& Cattafesta, L.~N. (2017).
De-biasing the dynamic mode decomposition for applied Koopman spectral
analysis of noisy datasets.
\textit{Theor.\ Comput.\ Fluid Dyn.}, 31(4), 349--368.

\bibitem[Jordan \& Dimanche(2021)]{jordan2021neural}
Jordan, M.~I. \& Dimanche, A. (2021).
\todo{Find exact reference for neural ODE periodic data instability.}

\bibitem[Kutz et~al.(2016)]{kutz2016dynamic}
Kutz, J.~N., Brunton, S.~L., Brunton, B.~W., \& Proctor, J.~L. (2016).
\textit{Dynamic Mode Decomposition: Data-Driven Modeling of Complex Systems}.
SIAM.

\bibitem[Kvalheim \& Revzen(2021)]{kvalheim2021existence}
Kvalheim, M.~D. \& Revzen, S. (2021).
Existence and uniqueness of global Koopman eigenfunctions for stable
fixed points and periodic orbits.
\textit{Physica D}, 425, 132959.

\bibitem[Li et~al.(2024)]{li2024adiabatic}
Li, M., Jain, S., \& Haller, G. (2024).
\todo{Exact adiabatic SSM reference.}

\bibitem[Linot et~al.(2023)]{linot2023stabilized}
Linot, A.~J., Burby, J.~W., Tang, Q., Balaprakash, P., Graham, M.~D.,
\& Maulik, R. (2023).
Stabilized neural ordinary differential equations for long-time
forecasting of dynamical systems.
\textit{J.\ Comput.\ Phys.}, 474, 111838.

\bibitem[Schmid(2010)]{schmid2010dynamic}
Schmid, P.~J. (2010).
Dynamic mode decomposition of numerical and experimental data.
\textit{J.\ Fluid Mech.}, 656, 5--28.

\bibitem[Strogatz(2015)]{strogatz2015nonlinear}
Strogatz, S.~H. (2015).
\textit{Nonlinear Dynamics and Chaos}.
Westview Press, 2nd edition.

\bibitem[Takens(1981)]{takens1981detecting}
Takens, F. (1981).
Detecting strange attractors in turbulence.
In \textit{Dynamical Systems and Turbulence}, pp.\ 366--381. Springer.

\bibitem[Trefethen \& Embree(2005)]{trefethen2005spectra}
Trefethen, L.~N. \& Embree, M. (2005).
\textit{Spectra and Pseudospectra: The Behavior of Nonnormal Matrices
and Operators}.
Princeton University Press.

\bibitem[Wilson \& Ermentrout(2016)]{wilson2016isostable}
Wilson, D. \& Ermentrout, B. (2016).
Greater accuracy and broadened applicability of phase reduction using
isostable coordinates.
\textit{J.\ Math.\ Biol.}, 76, 37--66.

\bibitem[Winner et~al.(2023)]{winner2023lstm}
Winner, T., \emph{et al.} (2023).
\todo{Full citation for Taniel's LSTM gait paper.}

\end{thebibliography}

\end{document}
